%================ Chapter ==================
\chapter{\ifenglish System Testing and Evaluation\else การทดสอบระบบ\fi}
\label{chap:eval}

%================ Section 4.1: Unit Testing ==================
\section{Unit Testing}
\label{sec:unit-testing}

\noindent
Unit Testing คือการทดสอบหน่วยย่อยที่สุดของโปรแกรมแบบแยกส่วน โดยแต่ละ test case ตรวจสอบว่า function หรือ method หนึ่งทำงานถูกต้องตาม specification ที่กำหนดไว้ภายใต้เงื่อนไขต่างๆ โดยไม่พึ่งพา component ภายนอก เช่น database หรือ network \cite{pressman2019}

GymMate เขียน unit test สำหรับ \textbf{Usecase Layer} ซึ่งเป็นชั้นที่บรรจุ business logic หลักของระบบ การที่ทดสอบได้โดยไม่ต้องพึ่ง database จริงเป็นผลโดยตรงจากการออกแบบตาม Clean Architecture ที่ Usecase Layer เข้าถึงข้อมูลผ่าน Repository interface เท่านั้น ทำให้สามารถแทนที่ด้วย mock object ระหว่างการทดสอบได้

\subsection{โครงสร้างของ Test}

แต่ละ feature module ภายใต้ \texttt{feature/\{feature\_name\}/usecase/} ประกอบด้วยไฟล์ 3 ไฟล์

\begin{itemize}
  \item \textbf{\texttt{usecase.go}} — implement business logic จริงของ feature รับ Repository interface ผ่าน constructor (Dependency Injection)
  \item \textbf{\texttt{mock\_test.go}} — สร้าง mock struct ที่ implement Repository interface เดียวกัน แต่ไม่ติดต่อ database จริง กำหนดพฤติกรรมได้ในแต่ละ test case
  \item \textbf{\texttt{usecase\_test.go}} — เขียน test case โดยสร้าง usecase instance ที่รับ mock repository แทน repository จริง แล้วตรวจสอบผลลัพธ์ที่ได้
\end{itemize}

\subsection{หลักการทดสอบ (Arrange-Act-Assert)}

test case แต่ละตัวออกแบบตามโครงสร้าง \textbf{Arrange-Act-Assert}

\begin{enumerate}
  \item \textbf{Arrange} — กำหนดพฤติกรรมของ mock และสร้าง input data
  \item \textbf{Act} — เรียก function ที่ต้องการทดสอบด้วย input ที่เตรียมไว้
  \item \textbf{Assert} — ตรวจสอบว่า output และ error ที่ได้ตรงตาม expected result
\end{enumerate}

แต่ละ function ทดสอบหลาย scenario ครอบคลุมทั้ง happy path และ error path

\begin{table}[h]
\centering
\begin{tabular}{@{}ll@{}}
\toprule
\textbf{Scenario} & \textbf{ตัวอย่าง (Equipment Usecase)} \\
\midrule
Happy path        & GetEquipment คืน equipment ที่ถูกต้องเมื่อ ID มีอยู่ \\
Not found         & GetEquipment คืน error เมื่อ ID ไม่มีในระบบ \\
Duplicate         & CreateEquipment คืน error เมื่อชื่อซ้ำ \\
Invalid input     & CreateEquipment คืน error เมื่อ input ไม่ผ่าน validation \\
Repository error  & คืน error ที่เหมาะสมเมื่อ database มีปัญหา \\
\bottomrule
\end{tabular}
\caption{ตัวอย่าง Test Scenarios ในระบบ GymMate}
\label{tab:test-scenarios}
\end{table}

\subsection{การรัน Test และ CI/CD}

test ทั้งหมดถูกรันอัตโนมัติใน \textbf{CI pipeline ผ่าน GitHub Actions} ทุกครั้งที่มีการ push code ไปยัง repository โดยใช้คำสั่ง \texttt{go test ./...}

หาก test ใดล้มเหลว CI pipeline จะหยุดทำงานและ block การ merge code เข้า main branch ทำให้มั่นใจได้ว่า business logic ที่ deploy ขึ้น production ผ่านการตรวจสอบแล้วทุกครั้ง

%================ Section 4.2: Usability Testing ==================
\section{การทดสอบการใช้งานจริง (Usability Testing / UAT)}
\label{sec:uat}

\noindent
การทดสอบนี้มีวัตถุประสงค์เพื่อประเมินประสบการณ์ผู้ใช้งาน (User Experience) และความง่ายในการใช้งานของระบบ

\subsection{กลุ่มผู้ทดสอบและรูปแบบการเก็บข้อมูล}

\begin{itemize}
  \item กลุ่มผู้ทดสอบ: นักศึกษาที่ใช้งานฟิตเนสและผู้ใช้งานทั่วไป จำนวน 11 คน
  \item รูปแบบการเก็บข้อมูล: แบบสอบถาม (Google Form) และการสังเกตพฤติกรรมผู้ใช้งาน
\end{itemize}

\subsection{หัวข้อการประเมิน}

\begin{itemize}
  \item ความง่ายในการใช้งาน (Ease of Use)
  \item ความเข้าใจของ UI
  \item ความสะดวกในการวางแผนการออกกำลังกาย
  \item ความพึงพอใจโดยรวม
  \item ความน่าจะเป็นในการใช้งานต่อ
  \item ความรู้สึกต่อดีไซน์ของแอป
\end{itemize}

\subsection{ผลการทดสอบ}

จากการเก็บข้อมูลผู้ทดสอบ 11 คน ผลการประเมินในภาพรวมแสดงให้เห็นว่า ผู้ทดสอบส่วนใหญ่พึงพอใจกับการออกแบบ UI/UX และมองว่าแอปใช้งานง่าย ฟีเจอร์ตอบโจทย์ความต้องการ และมีแนวโน้มจะใช้งานต่อในอนาคต

นอกจากนี้ในระหว่างการทดสอบ พบ bug 1 รายการ ได้แก่ การลืมใส่ Role middleware ที่ endpoint หนึ่ง ซึ่งได้รับการแก้ไขเรียบร้อยแล้ว

%================ Section 4.3: Objective Evaluation ==================
\section{การประเมินผลตามวัตถุประสงค์ของโครงการ}
\label{sec:objective-eval}

\noindent
จากวัตถุประสงค์ของโครงงานที่ตั้งไว้ในบทที่ 1 สามารถประเมินผลได้ดังตาราง \ref{tab:objective-eval}

\begin{table}[h]
\centering
\begin{tabular}{@{}p{5.5cm}p{7.5cm}@{}}
\toprule
\textbf{วัตถุประสงค์} & \textbf{ผลลัพธ์} \\
\midrule
พัฒนาระบบวางแผนการออกกำลังกาย (Workout Program Planner) &
  สามารถสร้างและใช้งานโปรแกรมทั้ง Preset Plan และ Custom Plan ได้จริง ผู้ใช้สามารถเลือก Goal \& Approach และดำเนินการออกกำลังกายตามปฏิทินได้ \\
\addlinespace
เพิ่มประสิทธิภาพด้วยระบบ Flex Substitute &
  ระบบสามารถแนะนำท่าทดแทนโดยใช้ Cosine Similarity เมื่ออุปกรณ์ไม่ว่างหรือผู้ใช้ต้องการเปลี่ยนท่า \\
\addlinespace
จัดการฟิตเนสผ่าน CMS Dashboard &
  Admin สามารถจัดการข้อมูลเครื่อง Floorplan Program และผู้ใช้ได้ผ่าน Web CMS \\
\addlinespace
สร้างแรงจูงใจในการออกกำลังกาย &
  ระบบ Gamification (XP, Badge, Streak, Leaderboard) ช่วยเพิ่ม engagement และการใช้งานอย่างต่อเนื่อง \\
\bottomrule
\end{tabular}
\caption{ผลการประเมินตามวัตถุประสงค์ของโครงงาน GymMate}
\label{tab:objective-eval}
\end{table}
